\section{The \dataname Benchmark}

\subsection{Overview and Design Philosophy}

\dataname is designed around a core hypothesis: safety evaluation in T2I systems should be gradient-aware rather than threshold-based.
Rather than providing a large flat set of binary-labeled examples, \dataname provides a smaller set of \emph{relational} examples (severity ladders) in which a single ladder of four prompts encodes more information about a model's safety boundary than four independent binary-labeled samples, because it exposes \emph{where} along the severity spectrum the model first fails.


This motivates our severity-graded design, in which every ladder in \dataname consists of four prompts and corresponding images, risk explanations escalating through four severity levels: Safe (L0), Low-risk (L1), Mid-risk (L2), and High-risk (L3). 
This structure enables analysis not just of whether a model fails at maximum risk, but of where along the severity spectrum its safety boundary lies and how sharply that boundary is defined. 
\red{To make the intermediate regime concrete, consider the per-rung risk explanations attached to each ladder: in the Violence ladder (\Fig{\ref{app:fig:sample_3}}), L1 depicts ``non-consensual physical assault with visible injury'' while L2 escalates to ``severe assault aftermath''; in the Self-Harm ladder (\Fig{\ref{app:fig:sample_4}}), L1 shows ``superficial marks, low-level harm with emotional suffering'' while L2 shows ``active self-harm with visible bleeding wounds''. Such explanations are provided for every ladder in Appendix~\ref{app:examples}, letting readers inspect directly what the L1--L2 boundary contains rather than inferring it from scores alone.}
% Models with a sharp, well-calibrated boundary are preferable to those that either fail at low severity (over-generation of harmful content) or succeed only at maximum severity (under-restriction that misses borderline cases).


\subsection{Hierarchical Risk Taxonomy}


We define a hierarchical risk taxonomy covering 11 top-level categories (\Fig{\ref{fig:overview}}), each representing a distinct dimension of potential harm in T2I generation. 
They were synthesized from major platform policies and existing literature, including open-source models DALL-E 3~\citep{dalle3}, Midjourney~\citep{midjourney_guideline}, LLaVAGuard~\citep{helff2024llavaguard}, T2ISafety~\citep{li2025t2isafety}, T2I-RiskyPrompt~\citep{zhang2025t2iriskyprompt}, along with proprietary models Google~\citep{google}, Microsoft~\citep{microsoft}, and Meta AI~\citep{metaAI}.
Each category is further organized into four severity levels with escalation rules (i.e., appearances at each risk category) (Appendix~\ref{app:taxonomy}, \Tab{\ref{app:tab:severity}}, \Tab{\ref{app:tab:taxonomy}}).



% calibrated to reflect escalating degrees of potential harm as shown in Appendix~\ref{app:taxonomy}, \Tab{\ref{app:tab:severity}}.
% We also define escalation rules (Appendix~\ref{app:taxonomy}, \Tab{\ref{app:tab:taxonomy}} (i.e., what should appear at each level of each risk category) to better instruct the data creation pipeline and evaluate the dataset. 
% Our policy covers both open-source and proprietary safety standards.
% To ensure structural consistency throughout the benchmark, the focus of these 11 categories is defined as follows: 


% We define a hierarchical risk taxonomy covering 11 top-level categories (see \Fig{\ref{fig:overview}}), each representing a distinct dimension of potential harm in T2I generation. 
% The categories were identified through a synthesis of existing T2I safety literature, content moderation guidelines from major platforms including open-source models DALL-E 3~\citep{dalle3}, Midjourney~\citep{midjourney_guideline}, LLaVAGuard~\citep{helff2024llavaguard}, T2ISafety~\citep{li2025t2isafety}, T2I-RiskyPrompt~\citep{zhang2025t2iriskyprompt}, along with proprietary models Google~\citep{google}, Microsoft~\citep{microsoft}, and Meta AI~\citep{metaAI}.
% Our policy is also reviewed by internal governance office.
% At a glance, the focus contents of 11 categories are:



% Each category is further organized into four severity levels, calibrated to reflect escalating degrees of potential harm as shown in Appendix~\ref{app:taxonomy}, \Tab{\ref{app:tab:severity}}.
% We also define escalation rules (Appendix~\ref{app:taxonomy}, \Tab{\ref{app:tab:taxonomy}} (i.e., what should appear at each level of each risk category) to better instruct the data creation pipeline and evaluate the dataset. 







% \subsection{Dataset Statistics}
% \dataname contains 1,026 ladders (4,104 prompts) spanning 11 categories; each rung holds a prompt, reference image, and risk explanation. 
% All ladders pass the monotonic quality filter; each category contributes 82 - 99 ladders (Appendix~\ref{app:statistics}, Table~\ref{app:tab:stats}). 
% Each rung's four prompts are independently generated (naturalness), and all four images share the same model and seed (generative variance control). \Tab{\ref{tab:comparison}} compares \dataname against existing ones.

\subsection{Dataset Statistics}
\dataname contains 1,026 ladders (4,104 prompts) spanning 11 categories; each rung holds a prompt, reference image, and risk explanation. All ladders pass the monotonic quality filter; each category contributes 82--99 ladders (Table~\ref{app:tab:stats}). Each rung's four prompts are independently generated (naturalness), and all four images share the same model and seed (generative variance control).
Categories are not mutually exclusive at high severity: cross-category co-occurrence analysis (Appendix~\ref{app:comatrix}) shows Violence--Disturbing as the highest overlap pair (18.9\%/21.4\%), while Hate Speech and Misinformation show moderate mutual overlap (11.8\%/12.6\%), reflecting fluid boundaries between semantically adjacent categories.
\Tab{\ref{tab:comparison}} compares \dataname against existing benchmarks.
\red{We note that the 1,026 ladders are the strict post-filtering \emph{yield} of the pipeline rather than an arbitrary quota: \dataname is designed to map continuous failure slopes at the L1/L2 boundary, an objective fundamentally distinct from the flat coverage targets of larger binary benchmarks; its compact size directly reflects multi-stage filtering optimized to eliminate label noise and ensure high data quality.}



% This gives $100$ scenarios per category $\times$ $11$ categories $\times$ $4$ severity levels = $4,400$ prompt-image pairs. 
% This results in $1,100$ unique scenarios ($100$ per category), providing a total of $4,400$ unique prompt-image pairs. 
% Each prompt-image pair is unique; prompts at each level are semantically related but independently generated to ensure naturalness and avoid template-based artifacts.


% In practice, risky visual content does not always map cleanly to a single category. 
% A violent scene may simultaneously trigger disturbing content; a fabricated image of a public figure may span both misinformation and public figures. 
% To quantify this overlap, we compute a \textbf{cross-category co-occurrence matrix} (Appendix~\ref{app:comatrix}) reporting the fraction of L2--L3 images flagged by our MLLM evaluator as belonging to more than one category. 
% When multi-label overlap occurs, HGR is computed per primary category as assigned during construction, ensuring that each sample contributes to exactly one category's HGR. 
% Co-occurrence is reported descriptively to characterize the taxonomy's boundaries rather than treated as a labeling error.


% SafeGrad contains 4,400 samples in total. Each sample corresponds to a unique risk scenario drawn from one of the 11 categories, instantiated at all four severity levels. This results in 1,100 unique scenarios (100 per category), providing a total of 4,400 unique prompt-image pairs. 


% \begin{table}[tb]
% \centering
% \caption{\dataname dataset statistics by category.}
% \label{tab:stats}
% \begin{tabular}{lrrrr}
% \toprule
% \textbf{Category} & \textbf{\# Scenarios} & \textbf{\# Prompts} & \textbf{\# Images} & \textbf{Avg. Prompt Length} \\
% \midrule
% Violence \& Physical Harm       & 100   & 400   & 400   & 18.3 words \\
% Sexual Content                  & 100   & 400   & 400   & 16.7 words \\
% Hate Speech \& Discrimination   & 100   & 400   & 400   & 19.1 words \\
% Self-Harm \& Suicide            & 100   & 400   & 400   & 17.4 words \\
% Dangerous \& Illegal Activities & 100   & 400   & 400   & 20.2 words \\
% Privacy Violations              & 100   & 400   & 400   & 15.8 words \\
% Political Extremism             & 100   & 400   & 400   & 21.0 words \\
% Misinformation \& Manipulation  & 100   & 400   & 400   & 18.9 words \\
% Copyright Infringement          & 100   & 400   & 400   & 14.3 words \\
% Stereotyping \& Bias            & 100   & 400   & 400   & 17.8 words \\
% Graphic Disturbing Content      & 100   & 400   & 400   & 16.2 words \\
% \midrule
% \textbf{Total}                  & \textbf{1,100} & \textbf{4,400} & \textbf{4,400} & \textbf{18.0 words} \\
% \bottomrule
% \end{tabular}
% \end{table}


\begin{table}[tb]
\centering
\resizebox{0.47\textwidth}{!}{%
\begin{tabular}{lrrrccc}
\toprule
\textbf{Dataset} & \textbf{\# Samples} & \textbf{\# Categories} & \textbf{\# Levels} & \textbf{Grad.} & \textbf{PPL ($\downarrow$)} & \textbf{Reason} \\
\midrule
UnsafeDiffusion~\citep{qu2023unsafe} & 434 & 5 & 2 & \xmark & 2,511  & \xmark \\
I2P~\citep{schramowski2023safe}               & 4,703  & 7  & 2 & \xmark & 2,587 & \xmark \\
T2ISafety~\citep{li2025t2isafety}           & ~70K & 12  & 2 & \xmark & 2,613 &  \xmark \\
T2I-RiskyPrompt~\citep{zhang2025t2iriskyprompt} & 6,432  & 14 & 2 & \xmark & 86 &  \cmark \\
\midrule
\textbf{SafeGrad (Ours)}                     & \textbf{4,104} & \textbf{11} & \textbf{4} & \cmark & 33.92 & \cmark  \\
\bottomrule
\end{tabular}%
}
\caption{Comparison of \dataname with existing benchmarks.
\textbf{Grad.} is severity-graded levels (vs.\ binary safe/unsafe); \textbf{PPL} is mean L3 prompt perplexity (lower means more fluent, so metaphorical high-risk prompts avoid explicit keywords and defeat perplexity-based anomaly detection simultaneously); \textbf{Reason.} is per-sample risk explanation annotation. \dataname is the only benchmark with graded severity, low-perplexity adversarial prompts, and risk explanations.}
\label{tab:comparison}
\end{table}
